\documentclass{cesi}

\cesisetup{
  author      = {Louis DROUHIN, Owen BOUDON, Sébastien RUET, Zakariya SAOULA},
  title       = {Breezy},
  subtitle    = {Réseau social distribué en architecture microservices},
  description = {Rapport de projet — Développement d'applications distribuées},
  typedoc     = {RAPPORT}
}

\begin{document}

\cesititlepage

\tableofcontents
\newpage

% =============================================================
\section{Introduction}
% =============================================================

\subsection{Contexte du module et du projet}

Ce rapport présente le travail réalisé dans le cadre du projet Fil Rouge du module
\emph{Développement d'applications distribuées}. Ce projet consiste à développer \textbf{Breezy},
un réseau social léger inspiré de Twitter/X.

Celui-ci se doit d'être robuste et scalable, en s'appuyant sur une architecture en
microservices plutôt que sur un monolithe applicatif.

\subsection{Objectifs et attentes du projet}

Le cahier des charges fourni en début de séquence définit Breezy comme un réseau
social inspiré de Twitter, dont le périmètre fonctionnel est découpé en
fonctionnalités primaires (obligatoires) et secondaires (optionnelles). Les
fonctionnalités primaires couvrent le socle minimal d'un réseau social viable :

\begin{itemize}
  \item création de compte et authentification sécurisée ;
  \item publication de messages courts (280 caractères) ;
  \item affichage des messages sur le profil de l'utilisateur ;
  \item fil d'actualité chronologique des comptes suivis ;
  \item interactions sociales de base : like, commentaire, réponse à un commentaire,
        follow/unfollow ;
  \item profil utilisateur avec informations de base.
\end{itemize}

\subsection{Présentation de l'équipe}

\begin{table}[H]
\centering
\begin{tabularx}{\linewidth}{l X}
\toprule
\textbf{Membre} & \textbf{Responsabilité principale} \\
\midrule
Louis DROUHIN     & \texttt{auth-svc} (authentification, JWT), \texttt{feed-svc}, et front-end \\
Sébastien RUET    & \texttt{user-svc} (profils, follow), gateway Nginx et front-end \\
Owen BOUDON       & \texttt{notif-svc} (notifications) \\
Zakariya SAOULA   & \texttt{post-svc} (publications, commentaires) \\
\bottomrule
\end{tabularx}
\caption{Répartition des responsabilités au sein de l'équipe}
\end{table}

Cette répartition suit le découpage naturel de l'architecture en microservices : à
partir du moment où les contrats d'API entre services ont été fixés, chaque membre a
pu avancer en parallèle sur son service sans dépendance bloquante sur le
travail des autres, à l'exception du front-end qui consomme l'ensemble des API.

\subsection{Périmètre retenu et hors-scope justifié}

Le cahier des charges propose un large éventail de fonctionnalités secondaires
(messagerie privée, médias riches, signalement de contenu, interface multi-langues,
thèmes personnalisés, suspension de comptes...). Conscients de la charge de travail
que représenterait leur implémentation complète au regard du temps imparti, nous avons
fait le choix de retirer volontairement plusieurs fonctionnalités du périmètre,
plutôt que de les implémenter partiellement ou de manière peu aboutie.

\begin{infobox}
\textbf{Fonctionnalités non développées, et raisons :}
\begin{enumerate}[label=\textbullet, leftmargin=1.4em, itemsep=4pt, topsep=6pt]
  \item \textbf{Vidéos sur les posts} : complexité de stockage et de traitement de
        fichiers non proportionnée au bénéfice pour cette itération,
        cependant les images et GIF sont pris en charge ;
  \item \textbf{Signalement de contenu et modération (suppression de posts, bannissements)} :
        nécessite un back-office supplémentaire et donc du temps supplémentaire,
        qui pourra cependant être développé dans un second temps ;
  \item \textbf{Thème sombre} : complexité de la conversion d'une palette de couleur en
        mode sombre, qui n'apporte pas de valeur fonctionnelle directe pour cette itération ;
  \item \textbf{Messagerie privée} : fonctionnalité complète à part entière
        (modèle de données, chiffrement éventuel) repoussée en évolution future.
\end{enumerate}
\end{infobox}

Ce choix de périmètre a été tranché collectivement et documenté dès le début du
développement, afin qu'aucun membre de l'équipe ne réintroduise une de ces
fonctionnalités par anticipation sans concertation.

\newpage
% =============================================================
\section{Analyse des besoins et fonctionnalités}
% =============================================================

\subsection{User stories et hiérarchisation}

Les fonctionnalités du cahier des charges sont exprimées sous forme de user stories
(« En tant que..., je veux..., afin de... »), ce qui a permis de prioriser le
développement par criticité plutôt que par ordre d'apparition dans le document.

\begin{table}[H]
\centering
\begin{tabularx}{\linewidth}{l X l}
\toprule
\textbf{Code} & \textbf{Fonctionnalité} & \textbf{Priorité} \\
\midrule
Fx1  & Création de compte avec validation        & Primaire \\
Fx2  & Authentification sécurisée (JWT)           & Primaire \\
Fx3  & Publication de messages courts             & Primaire \\
Fx4  & Affichage des messages sur le profil        & Primaire \\
Fx5  & Fil chronologique des comptes suivis        & Primaire \\
Fx6  & Like d'un post                              & Primaire \\
Fx7  & Commentaire sur un post                     & Primaire \\
Fx8  & Réponse à un commentaire                    & Primaire \\
Fx9  & Follow / unfollow                           & Primaire \\
Fx10 & Profil utilisateur (bio, avatar)             & Primaire \\
Fx11 & Liste des messages publiés sur le profil    & Primaire \\
Fx13 & Recherche par tags                          & Secondaire \\
Fx14--Fx16 & Notifications (mention, like, follower) & Secondaire \\
Fx18 & Ajout d'images aux messages                & Secondaire \\
Fx22 & Interface multi-langues                    & Secondaire \\
\bottomrule
\end{tabularx}
\caption{Hiérarchisation retenue des fonctionnalités}
\end{table}

\subsection{Matrice des permissions}

La gestion des droits repose sur un rôle porté par chaque compte (\texttt{user} ou
\texttt{admin}), vérifié à la gateway. Le tableau suivant synthétise les accès
retenus pour les fonctionnalités implémentées.

\begin{table}[H]
\centering
\begin{tabularx}{\linewidth}{X c c}
\toprule
\textbf{Fonctionnalité} & \textbf{Visiteur} & \textbf{Utilisateur authentifié} \\
\midrule
Création de compte               & Oui & --- \\
Authentification                  & Oui & Oui \\
Publication de messages           & Non & Oui \\
Consultation des profils/posts    & Non & Oui \\
Like, commentaire, réponse        & Non & Oui \\
Follow / unfollow                 & Non & Oui \\
\bottomrule
\end{tabularx}
\caption{Matrice des permissions appliquée au périmètre retenu}
\end{table}

\subsection{Méthodologie et gestion de projet}

L'équipe s'est appuyée sur des standards tout au long du développement :

\begin{itemize}
  \item un modèle de branches \textbf{Gitflow} (\texttt{main}, \texttt{develop},
        \texttt{feature/<service>-<description>}, \texttt{fix/<description>}),
        chaque fonctionnalité étant développée à partir de \texttt{develop} et jamais
        directement sur \texttt{main} ;
  \item des messages de commit normalisés selon la convention
        \textbf{Conventional Commits} (\texttt{feat}, \texttt{fix}, \texttt{docs},
        \texttt{refactor}, \texttt{test}, \texttt{chore}), avec le service concerné en
        \emph{scope} (ex. \texttt{feat(auth-svc): ajoute la route de login}) ;
  \item un \textbf{board Kanban} à trois colonnes principales
        (À faire / En cours / Terminé), complété d'une colonne Blocage pour signaler
        les dépendances entre services ;
  \item un \textbf{dépôt privé} partagé entre les quatre membres.
\end{itemize}

\subsection{Détail des étapes et ressources mobilisées}

Le projet a suivi la progression du module : modélisation de l'architecture
(boucle API), développement du back-end REST par service, intégration de la
sécurité JWT et de la conteneurisation Docker, puis développement du front-end
React/Next.js consommant les API exposées par la gateway.

\begin{table}[H]
\centering
\begin{tabularx}{\linewidth}{l X}
\toprule
\textbf{Catégorie} & \textbf{Ressources mobilisées} \\
\midrule
Langage / runtime   & JavaScript / TypeScript, Node.js \\
Back-end            & Express.js, Sequelize, Mongoose, jsonwebtoken \\
Bases de données    & PostgreSQL (\texttt{auth-svc}, \texttt{user-svc}),
                       MongoDB (\texttt{post-svc}, \texttt{notif-svc}) \\
Front-end            & React, Next.js \\
Infrastructure       & Docker, Docker Compose, Nginx (reverse proxy) \\
Outillage            & pnpm (monorepo), Git, GitHub (dépôt, Issues, Projects) \\
\bottomrule
\end{tabularx}
\caption{Ressources techniques mobilisées pour le projet}
\end{table}

\newpage
% =============================================================
\section{Architecture du système}
% =============================================================

\subsection{Vue d'ensemble et justification du choix microservices}

Le sujet imposait de réfléchir au choix d'une architecture logicielle adaptée,
entre monolithique et microservices. Nous avons retenu une architecture en
microservices, chaque service disposant de sa propre base de données, 
pour les raisons suivantes :

\begin{itemize}
  \item \textbf{isolation des responsabilités} : authentification, gestion de profils,
        publications et notifications sont des domaines fonctionnels distincts, avec
        des contraintes de cohérence et de volumétrie différentes ;
  \item \textbf{parallélisation du développement} : le découpage par service a permis
        à chaque membre de l'équipe de travailler de façon autonome dès que les
        contrats d'API ont été stabilisés ;
  \item \textbf{portabilité} : chaque service est conteneurisé indépendamment, ce qui
        facilite un futur déploiement sur un orchestrateur (cf. axes d'amélioration).
\end{itemize}

\subsection{Schéma d'architecture}

\begin{infobox}
\textbf{Architecture retenue} : un point d'entrée HTTP unique (Nginx) route les
requêtes vers cinq services métier, isolés du frontend sur leur propre réseau docker
\texttt{backend} et disposant (sauf \texttt{feed-svc}) de sa propre base de données.
\end{infobox}

\begin{table}[H]
\centering
\begin{tabularx}{\linewidth}{l l l}
\toprule
\textbf{Service} & \textbf{Rôle} & \textbf{Port par défaut} \\
\midrule
\texttt{auth-svc}  & Authentification, JWT             & 3001 \\
\texttt{user-svc}  & Profils, relations follow         & 3002 \\
\texttt{post-svc}  & Publications, commentaires        & 3003 \\
\texttt{notif-svc} & Notifications                     & 3004 \\
\texttt{feed-svc}  & Fil d'actualité (sans base)        & 3005 \\
Frontend Next.js   & Interface utilisateur             & 3000 \\
Nginx (gateway)    & Point d'entrée, reverse proxy     & 80 \\
\bottomrule
\end{tabularx}
\caption{Services et ports de l'architecture Breezy}
\end{table}

Le client ne dialogue jamais directement avec un service métier : toutes les
requêtes transitent par Nginx, qui les redirige selon le préfixe de route
(\texttt{/api/auth}, \texttt{/api/users}, \texttt{/api/posts},
\texttt{/api/notifications}, \texttt{/api/feed}). Aucun service backend n'est
exposé directement à l'extérieur du réseau Docker.

\subsection{Stack technique par service}

\begin{table}[H]
\centering
\begin{tabularx}{\linewidth}{l l l X}
\toprule
\textbf{Service} & \textbf{Moteur} & \textbf{ORM/ODM} & \textbf{Raison du choix} \\
\midrule
\texttt{auth-svc}  & PostgreSQL & Sequelize & Données critiques, contraintes fortes,
                                              transactions ACID \\
\texttt{user-svc}  & PostgreSQL & Sequelize & Le graphe de follow est un
                                              many-to-many relationnel classique \\
\texttt{post-svc}  & MongoDB    & Mongoose  & Volume d'écriture élevé, structure
                                              semi-flexible \\
\texttt{notif-svc} & MongoDB    & Mongoose  & Fort volume d'écriture, payload
                                              variable selon le type de notification \\
\texttt{feed-svc}  & ---        & ---       & Calculé à la volée, pas de
                                              persistance propre \\
\bottomrule
\end{tabularx}
\caption{Choix technologiques par service et justification}
\end{table}

\subsection{Modèle de données et cohérence inter-services}

Le \texttt{username} constitue la clé de cohérence partagée entre tous les
services :

\begin{itemize}
  \item Il est immuable : aucune route de modification de username n'existe,
        dans aucun service ;
  \item Il est la clé primaire des entités utilisateur (\texttt{Account}
        dans \texttt{auth-svc}, \texttt{Profile} dans \texttt{user-svc}) ;
  \item Les services métier (\texttt{post-svc}, \texttt{notif-svc}) ne font aucune
        clé étrangère technique vers \texttt{auth-svc}/\texttt{user-svc} (impossible
        de toute façon entre bases de données différentes) mais stockent simplement
        le \texttt{username} en tant que champ simple
        (\texttt{authorUsername}, \texttt{recipientUsername}) ;
  \item La cohérence est donc applicative et non technique : c'est
        \texttt{auth-svc} qui garantit l'unicité du username à la création de compte,
        les autres services lui font confiance.
\end{itemize}

Un choix de modélisation notable concerne les commentaires dans \texttt{post-svc} :
un commentaire est un \texttt{Post} comme un autre, simplement doté d'un
champ \texttt{parentId} pointant vers le post auquel il répond. Ce choix évite de
dupliquer un modèle \texttt{Comment} séparé et permet nativement des fils de
discussion à profondeur illimitée.

\subsection{Sécurité et gateway Nginx}

Conformément aux standards vus dans le prosit JWT/Docker, l'authentification
repose sur un mécanisme JWT délégué :

\begin{enumerate}
  \item Nginx intercepte chaque requête entrante et déclenche une directive
        \texttt{auth\_request} ;
  \item cette directive délègue la validation effective du jeton à
        \texttt{auth-svc}, via une route dédiée (\texttt{/api/auth/validate}) ;
  \item si le jeton est valide, Nginx transmet aux services métier deux en-têtes
        de confiance : \texttt{X-User-Username} et \texttt{X-User-Role} ;
  \item les services métier (\texttt{post-svc}, \texttt{notif-svc}, \texttt{user-svc},
        \texttt{feed-svc}) ne réimplémentent pas de logique de vérification JWT : ils
        font confiance aux en-têtes transmis par la gateway après validation.
\end{enumerate}

\begin{infobox}
\textbf{Gateway = Nginx pur.} Aucun service Node.js personnalisé ne fait office
d'API Gateway dans ce projet : Nginx, configuré en reverse proxy, assure à lui seul
le routage et la délégation de l'authentification. Ce choix limite la surface de
code à maintenir et s'appuie sur un composant éprouvé plutôt que de réinventer une
brique de sécurité.
\end{infobox}

Les autres règles de sécurité transverses appliquées à l'ensemble des services
Node.js sont : gestion d'erreurs centralisée (middleware d'erreur Express, jamais
de stack trace brute renvoyée au client), CORS configuré explicitement, variables
sensibles exclusivement via \texttt{dotenv}, et respect du niveau 2 du modèle de
maturité de Richardson (verbes HTTP sémantiquement corrects, codes de statut HTTP).

\subsection{Conteneurisation}

L'ensemble des services est orchestré via Docker Compose, avec un fichier dédié au
développement (hot-reload par bind-mount du code source) et un fichier dédié à la
production (\texttt{docker-compose.prod.yml}, images distinctes par microservice,
construites et hébérgées sur le Docker Registry de Github via Github Actions).
Deux réseaux Docker isolent le frontend du backend, Nginx faisant le pont entre les
deux. Des \emph{profiles} Docker Compose permettent de ne démarrer que les services
utiles pour du développment ciblé plutôt que l'intégralité de la stack.

\newpage
% =============================================================
\section{Réalisation}
% =============================================================

Cette section détaille, service par service, les fonctionnalités implémentées et
les choix de mise en œuvre. Chaque sous-section a été rédigée par le ou les membres
ayant porté le développement du service correspondant.

\subsection{\texttt{auth-svc} : Authentification et JWT (Louis DROUHIN)}

\texttt{auth-svc} couvre Fx1 (création de compte) et Fx2 (authentification sécurisée).
Il ne stocke que les informations strictement nécessaires à l'authentification :
\texttt{username}, \texttt{email}, \texttt{passwordHash}, \texttt{role},
\texttt{active}. Toute donnée de profil (bio, avatar) relève de \texttt{user-svc}.

\begin{itemize}
  \item inscription avec hachage du mot de passe (jamais stocké en clair) ;
  \item connexion délivrant un jeton JWT signé, vérifié ensuite côté Nginx via
        \texttt{auth\_request} ;
  \item déconnexion par suppression de la ligne de refresh token en base, sans champ
        de révocation dédié, choix retenu pour limiter la complexité du modèle de
        données au regard du périmètre du projet.
\end{itemize}

\subsection{\texttt{user-svc} : Profils utilisateurs (Sébastien RUET)}

\texttt{user-svc} couvre Fx9 (follow/unfollow) et Fx10 (profil utilisateur). Le
graphe de relations de follow est modélisé de façon relationnelle classique
(table d'association many-to-many), ce qui justifie le choix de PostgreSQL plutôt
que MongoDB pour ce service.

\subsection{\texttt{post-svc} : Publications et commentaires (Zakariya SAOULA)}

\texttt{post-svc} couvre Fx3 (publication de messages courts), Fx4 (affichage sur le
profil), Fx6 (like), Fx7 et Fx8 (commentaires et réponses). Le choix de MongoDB
pour ce service est justifié par un volume d'écriture potentiellement élevé et une
structure de données semi-flexible (un post peut porter un nombre variable de
likes, de tags, de réponses).

Le point de mise en œuvre central est le modèle unique \texttt{Post} avec champ
\texttt{parentId} optionnel : un post racine a un \texttt{parentId} nul, un
commentaire référence le post auquel il répond, et une réponse à un commentaire
référence ce commentaire, sans limite de profondeur. Le compteur de réponses
(\texttt{replyCount}) et de like (\texttt{likeCount}) est recalculé à chaque ajout/suppression
pour éviter un comptage à la volée coûteux sur des fils profonds.

\subsection{\texttt{notif-svc} : Notifications (Owen BOUDON)}

\texttt{notif-svc} couvre Fx14 à Fx16 (notifications de mention, de like, de nouveau
follower). Le choix de MongoDB se justifie par un fort volume d'écriture et un
\emph{payload} variable selon le type de notification (une notification de mention
ne porte pas les mêmes champs qu'une notification de like).

Le service expose une route de consultation des notifications d'un utilisateur
(authentifié via les en-têtes transmis par la gateway) ainsi qu'une route de
marquage comme lues, sans dépendre directement des bases de données des autres
services : la création d'une notification est déclenchée applicativement par
\texttt{post-svc} ou \texttt{user-svc} au moment de l'action concernée.

\subsection{\texttt{feed-svc} : Fil d'actualité (Louis DROUHIN)}

\texttt{feed-svc} couvre Fx5 (fil chronologique des comptes suivis). C'est le seul
service de l'architecture sans base de données propre : le fil est calculé
à la volée (\emph{fan-out on read}). À chaque consultation :

\begin{enumerate}
  \item \texttt{feed-svc} interroge \texttt{user-svc} pour obtenir la liste des
        comptes suivis par l'utilisateur courant ;
  \item il interroge ensuite \texttt{post-svc} pour récupérer les posts de ces
        comptes ;
  \item les résultats sont triés par date et renvoyés au client.
\end{enumerate}

\begin{infobox}
Ce choix de conception est documenté comme un compromis assumé plutôt qu'une
contrainte technique : le \emph{fan-out on read} est plus simple à mettre en œuvre
qu'un fil matérialisé, au prix d'un coût de lecture plus élevé. Il est cohérent
avec le volume de données attendu dans le cadre du projet, et identifié comme axe
d'amélioration en cas de montée en charge (cf. section 6).

De plus, le fait qu'il soit
dans un service dédié permet d'envisager un développement futur d'un algorithme
de recommendation.
\end{infobox}

\subsection{Gateway Nginx (Sébastien RUET)}

La configuration de la gateway Nginx (\texttt{gateway/nginx.conf})
définit le routage vers chaque service métier ainsi que la directive
\texttt{auth\_request} déléguant la validation JWT à \texttt{auth-svc}, conformément
au choix d'architecture présenté en section 3.5.

\subsection{Front-end (Louis DROUHIN, Sébastien RUET)}

Le front-end est développé en React/Next.js, avec une approche mobile-first. Il
consomme exclusivement les API exposées via la gateway Nginx (jamais directement
les services métier) et porte les écrans suivants : authentification (connexion,
inscription), fil d'actualité, profil utilisateur, et les interactions sur les
posts (publication, like, commentaire, réponse).

\newpage
% =============================================================
\section{Difficultés rencontrées et solutions}
% =============================================================

\subsection{Choix techniques discutés en équipe}

\begin{itemize}
  \item \textbf{Fan-out on read vs fan-out on write} pour le fil d'actualité : le
        choix du calcul à la volée a été préféré à un fil matérialisé pour limiter
        la complexité du modèle de données dans le temps imparti, au prix d'une
        charge de lecture plus élevée à chaque consultation du fil.
  \item \textbf{Modélisation des commentaires} : la tentation initiale de créer une
        entité \texttt{Comment} séparée de \texttt{Post} aurait limité la profondeur
        des fils de discussion à un seul niveau ; le choix d'unifier les deux entités
        sous un \texttt{Post} avec \texttt{parentId} a permis de lever cette
        limitation sans complexifier le modèle de données.
  \item \textbf{Délégation de la sécurité à la gateway} : centraliser la validation
        JWT dans Nginx plutôt que de la dupliquer dans chaque service métier a
        nécessité de figer tôt le contrat des en-têtes transmis
        (\texttt{X-User-Username}, \texttt{X-User-Role}) pour que tous les services
        s'alignent sans ambiguïté.
\end{itemize}

\subsection{Points de friction inter-services}

Le développement en parallèle de cinq services par des membres différents a
nécessité de fixer rapidement des conventions communes (arborescence de dossiers,
format des routes, gestion des erreurs) afin d'éviter une divergence de pratiques
entre services qui aurait complexifié l'intégration finale via la gateway. Le
\texttt{username} comme clé de cohérence partagée entre services a constitué le
principal point d'attention, en l'absence de clé étrangère technique possible entre
des bases de données distinctes.

\newpage
% =============================================================
\section{Conclusion}
% =============================================================

\subsection{Bilan par rapport aux objectifs}

Le projet Breezy a permis de mettre en application l'ensemble des compétences
travaillées au long de la séquence : conception argumentée d'une architecture
microservices, développement d'API REST de niveau 2 de maturité, sécurisation par
JWT déléguée à une gateway, développement d'un front-end React mobile-first, et
conteneurisation complète de l'environnement. L'ensemble des fonctionnalités
primaires du cahier des charges (Fx1 à Fx11) a été couvert, dans le périmètre
explicitement retenu par l'équipe.

\subsection{Axes d'amélioration et évolutions potentielles}

\begin{itemize}
  \item \textbf{Fan-out on write} pour le fil d'actualité, si la charge de lecture
        devenait significative, en matérialisant le fil de chaque utilisateur à
        l'écriture plutôt qu'à la consultation ;
  \item \textbf{Messagerie privée} (Fx17) et \textbf{signalement de contenu} (Fx20),
        retirés du périmètre initial, pour renforcer la dimension communautaire et
        la modération ;
  \item \textbf{Vidéos} sur les posts (Fx19) ;
  \item \textbf{Révocation explicite des jetons} (liste de révocation) plutôt que la
        simple suppression de la ligne en base, pour une gestion plus fine des
        sessions multiples ;
  \item \textbf{Migration vers un orchestrateur} (Kubernetes ou équivalent managé),
        dans la continuité de la portabilité déjà recherchée par la conteneurisation
        Docker, pour anticiper une montée en charge ;
  \item \textbf{Tests automatisés} (unitaires et d'intégration) sur chaque service,
        non couverts dans le temps imparti à cette itération.
\end{itemize}

\end{document}
